\documentclass[12pt,a4paper]{amsart}
\usepackage{amsmath,amssymb,amsthm}
\usepackage{mathtools}
\usepackage[utf8]{inputenc}
\usepackage[T1]{fontenc}
\usepackage[ngerman]{babel}
\usepackage{hyperref}
\usepackage{enumitem,booktabs,array}
\usepackage{geometry}
\geometry{margin=2.5cm}

\newtheorem{theorem}{Satz}[section]
\newtheorem{lemma}[theorem]{Lemma}
\newtheorem{corollary}[theorem]{Korollar}
\newtheorem{proposition}[theorem]{Proposition}
\newtheorem{conjecture}[theorem]{Vermutung}
\theoremstyle{definition}
\newtheorem{definition}[theorem]{Definition}
\newtheorem{example}[theorem]{Beispiel}
\newtheorem{remark}[theorem]{Bemerkung}

\author{Michael Fuhrmann}
\address{Specialist Mathematics Project, Build 141}
\email{specialist-maths@localhost}
\date{\today}

\begin{document}
\title{Spezielle Funktionen der Mathematik und Physik:\\
Gamma-, Bessel-, hypergeometrische und verwandte Funktionen}

\begin{abstract}
Spezielle Funktionen entstehen als kanonische Lösungen der fundamentalsten Differentialgleichungen
der mathematischen Physik und Analysis.  Dieses Paper entwickelt systematisch die wichtigsten
speziellen Funktionen: die Gamma- und Beta-Funktion, Bessel-Funktionen erster und zweiter Art,
die Gaußsche hypergeometrische Funktion \({}_2F_1\), Legendre-Polynome und zugeordnete
Legendre-Funktionen, Hermite- und Laguerre-Polynome mit ihren quantenmechanischen Anwendungen,
elliptische Funktionen von Weierstra\ss{} und Jacobi sowie die Riemann-Zetafunktion mitsamt
Dirichlet-\(L\)-Funktionen.  Für jede Funktionenklasse werden die definierende
Differentialgleichung bzw.\ Reihe, fundamentale Identitäten wie die Reflexionsformel
\(\Gamma(z)\Gamma(1-z)=\pi/\sin(\pi z)\), die Beta-Gamma-Relation und die
Funktionalgleichung von \(\zeta(s)\) bewiesen sowie die Verbindungen zwischen den
verschiedenen Klassen hergestellt.  Numerische Verifikationen aus dem begleitenden
Python-Modul \texttt{special\_functions.py} bestätigen alle theoretischen Ergebnisse.
\end{abstract}

\maketitle

\tableofcontents

%% ============================================================
\section{Einleitung}
%% ============================================================

Spezielle Funktionen sind diejenigen Funktionen, die so häufig bei der Lösung kanonischer
Probleme in Analysis, Geometrie und mathematischer Physik auftreten, dass sie eigene Namen
und umfangreiche Tabellen ihrer Eigenschaften erhalten haben.  Ihre systematische Untersuchung
zählt zu den schönsten Kapiteln der klassischen Analysis.

Die zentralen Objekte dieses Papers sind Lösungen linearer gewöhnlicher Differentialgleichungen
zweiter Ordnung mit Polynomkoeffizienten.  Nach dem Satz von Fuchs besitzen solche
Gleichungen in der Umgebung regulärer singulärer Punkte Potenzreihenlösungen; die so
entstehenden Funktionen weisen eine reiche algebraische und analytische Struktur auf:
Rekurrenzrelationen, erzeugende Funktionen, Orthogonalität auf einem geeigneten Intervall,
asymptotische Entwicklungen und Verbindungen zur Zahlentheorie.

\begin{example}[Kanonische Differentialgleichungen]
Die folgenden Differentialgleichungen zweiter Ordnung definieren je eine fundamentale
Familie spezieller Funktionen:
\begin{align}
  x^2 y'' + x y' + (x^2 - \nu^2)y &= 0 \qquad \text{(Bessel-Gleichung)},\label{eq:bessel_ode}\\
  (1-x^2)y'' - 2xy' + n(n+1)y &= 0 \qquad \text{(Legendre-Gleichung)},\label{eq:legendre_ode}\\
  z(1-z)w'' + [c-(a+b+1)z]w' - ab\,w &= 0 \qquad \text{(hypergeometrische Gleichung)},\label{eq:hyp_ode}\\
  y'' &= x\,y \qquad \text{(Airy-Gleichung)}.\label{eq:airy_ode}
\end{align}
\end{example}

Das Paper ist wie folgt aufgebaut.
Abschnitt~\ref{sec:gamma} behandelt die Gamma-Funktion einschließlich Reflexions- und
Verdopplungsformel.  Abschnitt~\ref{sec:beta} stellt die Beta-Funktion und ihre
Integraldarstellung vor.  Abschnitt~\ref{sec:bessel} entwickelt Bessel-Funktionen.
Abschnitt~\ref{sec:hypergeometric} behandelt die hypergeometrische Funktion \({}_2F_1\).
Abschnitt~\ref{sec:legendre} präsentiert Legendre-Polynome und Kugelflächenfunktionen.
Abschnitt~\ref{sec:hermite_laguerre} deckt Hermite- und Laguerre-Polynome sowie ihre
Rolle in der Quantenmechanik ab.  Abschnitt~\ref{sec:elliptic} behandelt elliptische
Funktionen.  Abschnitt~\ref{sec:zeta} diskutiert die Riemannsche Zetafunktion und
Dirichlet-\(L\)-Funktionen.  Abschnitt~\ref{sec:anwendungen} gibt einen Überblick über
physikalische und zahlentheoretische Anwendungen.

%% ============================================================
\section{Die Gamma-Funktion}
\label{sec:gamma}
%% ============================================================

\subsection{Definition und grundlegende Eigenschaften}

\begin{definition}[Gamma-Funktion via Integral]
Für \(\Re(z) > 0\) setze man
\[
  \Gamma(z) \;=\; \int_0^\infty t^{z-1} e^{-t}\,dt.
\]
\end{definition}

Partielle Integration liefert sofort die Funktionalgleichung.

\begin{theorem}[Funktionalgleichung]\label{thm:gamma_functional_de}
Für alle \(z\) mit \(\Re(z) > 0\) gilt
\[
  \Gamma(z+1) = z\,\Gamma(z).
\]
Insbesondere ist \(\Gamma(n+1) = n!\) für jedes \(n \in \mathbb{N}_0\).
\end{theorem}

\begin{proof}
Partielle Integration ergibt:
\[
  \Gamma(z+1)
  = \int_0^\infty t^z e^{-t}\,dt
  = \bigl[-t^z e^{-t}\bigr]_0^\infty + z\int_0^\infty t^{z-1}e^{-t}\,dt
  = 0 + z\,\Gamma(z). \qquad\qquad\square
\]
\end{proof}

Die Funktionalgleichung zusammen mit \(\Gamma(1) = 1\) liefert die meromorphe
Fortsetzung von \(\Gamma\) auf ganz \(\mathbb{C}\): \(\Gamma\) hat einfache Pole in
\(z = 0, -1, -2, \ldots\) und keine Nullstellen.

\subsection{Weierstraßsches Produkt und Gaußsche Formel}

Die Euler-Mascheroni-Konstante \(\gamma = -\Gamma'(1) \approx 0{,}5772\) tritt im
Weierstraßschen Produktansatz auf:
\[
  \frac{1}{\Gamma(z)}
  = z\,e^{\gamma z}\prod_{n=1}^\infty \left(1+\frac{z}{n}\right)e^{-z/n}.
\]
Äquivalent dazu ist die Gaußsche Grenzwertformel
\[
  \Gamma(z) = \lim_{n\to\infty} \frac{n!\,n^z}{z(z+1)\cdots(z+n)}.
\]

\subsection{Reflexionsformel}

\begin{theorem}[Eulersche Reflexionsformel]\label{thm:gamma_reflection_de}
Für alle \(z \notin \mathbb{Z}\) gilt
\[
  \Gamma(z)\,\Gamma(1-z) = \frac{\pi}{\sin(\pi z)}.
\]
\end{theorem}

\begin{proof}
Aus dem Weierstraßschen Produkt für \(1/\Gamma(z)\) und \(1/\Gamma(-z)\) folgt
\[
  \frac{1}{\Gamma(z)\,\Gamma(-z)}
  = -z\prod_{n=1}^\infty\left(1-\frac{z^2}{n^2}\right).
\]
Das Eulersche Produkt für den Sinus liefert
\[
  \frac{\sin(\pi z)}{\pi z} = \prod_{n=1}^\infty\left(1-\frac{z^2}{n^2}\right),
\]
woraus \(\Gamma(z)\Gamma(-z) = -\pi/(z\sin(\pi z))\) folgt.
Die Substitution \(z\mapsto 1-z\) und \(\Gamma(1-z) = -z\,\Gamma(-z)\) liefern
das Ergebnis.
\end{proof}

\begin{corollary}
\(\Gamma\!\left(\tfrac{1}{2}\right) = \sqrt{\pi}\).
\end{corollary}
\begin{proof}
Setzt man \(z = 1/2\) in Satz~\ref{thm:gamma_reflection_de}: \(\Gamma(1/2)^2 = \pi/\sin(\pi/2) = \pi\).
\end{proof}

\subsection{Verdopplungsformel}

\begin{theorem}[Legendresche Verdopplungsformel]\label{thm:gamma_duplication_de}
Für alle \(z\) mit \(\Re(z) > 0\) gilt
\[
  \Gamma(z)\,\Gamma\!\left(z+\tfrac{1}{2}\right)
  = \frac{\sqrt{\pi}}{2^{2z-1}}\,\Gamma(2z).
\]
\end{theorem}

\begin{proof}
Betrachte das Beta-Integral \(B(z,z) = \Gamma(z)^2/\Gamma(2z)\).  Die Substitution
\(t = (1+u)/2\) führt auf
\[
  B(z,z) = 2^{1-2z}\int_{-1}^1 (1-u^2)^{z-1}\,du = 2^{1-2z}\,B\!\left(\tfrac{1}{2},z\right)
  = 2^{1-2z}\frac{\sqrt{\pi}\,\Gamma(z)}{\Gamma(z+1/2)}.
\]
Umstellen liefert die behauptete Formel.
\end{proof}

\subsection{Stirlingsches Asymptotik}

Für großes \(|z|\) mit \(|\arg z| < \pi\) gilt
\[
  \ln\Gamma(z) = \left(z-\tfrac{1}{2}\right)\ln z - z + \tfrac{1}{2}\ln(2\pi)
  + \sum_{k=1}^N \frac{B_{2k}}{2k(2k-1)z^{2k-1}} + O\!\left(z^{-2N-1}\right),
\]
wobei \(B_{2k}\) die Bernoulli-Zahlen sind.  Insbesondere gilt die Stirlingsche Formel
\[
  n! \;\sim\; \sqrt{2\pi n}\,\left(\frac{n}{e}\right)^n \quad (n\to\infty).
\]

%% ============================================================
\section{Die Beta-Funktion}
\label{sec:beta}
%% ============================================================

\begin{definition}[Beta-Funktion]
Für \(\Re(a) > 0\), \(\Re(b) > 0\) sei
\[
  B(a,b) = \int_0^1 t^{a-1}(1-t)^{b-1}\,dt.
\]
\end{definition}

\begin{theorem}[Beta-Gamma-Relation]\label{thm:beta_gamma_de}
\[
  B(a,b) = \frac{\Gamma(a)\,\Gamma(b)}{\Gamma(a+b)}.
\]
\end{theorem}

\begin{proof}
Man schreibe \(\Gamma(a)\Gamma(b)\) als Doppelintegral:
\[
  \Gamma(a)\Gamma(b)
  = \int_0^\infty s^{a-1}e^{-s}\,ds\int_0^\infty t^{b-1}e^{-t}\,dt.
\]
Die Substitution \(s = r\,u\), \(t = r(1-u)\) mit \(r \in (0,\infty)\), \(u \in (0,1)\)
und Jakobi-Determinante \(r\) ergibt
\[
  = \int_0^\infty r^{a+b-1}e^{-r}\,dr \cdot \int_0^1 u^{a-1}(1-u)^{b-1}\,du
  = \Gamma(a+b)\,B(a,b). \qquad\square
\]
\end{proof}

\begin{proposition}[Symmetrie]
\(B(a,b) = B(b,a)\).
\end{proposition}

\begin{proposition}[Trigonometrische Form]
\[
  B(a,b) = 2\int_0^{\pi/2}\sin^{2a-1}\theta\,\cos^{2b-1}\theta\,d\theta.
\]
\end{proposition}

\begin{remark}
Die Beta-Funktion erscheint natürlich bei der Berechnung von Volumina der \(n\)-dimensionalen
Kugel, in der Theorie der Dirichlet-Verteilung und als Normierungskonstante orthogonaler
Polynome.
\end{remark}

%% ============================================================
\section{Bessel-Funktionen}
\label{sec:bessel}
%% ============================================================

\subsection{Die Bessel-Differentialgleichung}

\begin{definition}[Bessel-Gleichung der Ordnung \(\nu\)]
\begin{equation}\label{eq:bessel_de}
  x^2 y'' + x y' + (x^2 - \nu^2)y = 0, \qquad x > 0,\; \nu \in \mathbb{C}.
\end{equation}
\end{definition}

Gleichung~\eqref{eq:bessel_de} besitzt bei \(x = 0\) einen regulären singulären Punkt
mit Indizialwurzeln \(\pm\nu\).  Die Frobenius-Methode liefert zwei linear unabhängige
Lösungen \(J_\nu(x)\) und \(Y_\nu(x)\) (für nicht-ganzzahlige \(\nu\); für ganzzahlige
\(\nu\) enthält \(Y_\nu\) einen Logarithmusterm).

\subsection{Bessel-Funktionen erster Art}

\begin{definition}[Bessel-Funktion \(J_\nu\)]
\[
  J_\nu(x) = \sum_{m=0}^\infty \frac{(-1)^m}{m!\,\Gamma(m+\nu+1)}
  \left(\frac{x}{2}\right)^{2m+\nu}.
\]
\end{definition}

Für ganzzahlige \(n \in \mathbb{Z}\) gilt die Integraldarstellung
\[
  J_n(x) = \frac{1}{\pi}\int_0^\pi \cos(n\tau - x\sin\tau)\,d\tau.
\]

\subsection{Bessel-Funktionen zweiter Art (Neumann-Funktionen)}

\begin{definition}[Neumann-Funktion \(Y_\nu\)]
Für nicht-ganzzahlige \(\nu\) sei
\[
  Y_\nu(x) = \frac{J_\nu(x)\cos(\nu\pi) - J_{-\nu}(x)}{\sin(\nu\pi)}.
\]
Für ganzzahlige \(n\) definiert man \(Y_n(x) = \lim_{\nu\to n} Y_\nu(x)\), was
einen Logarithmusterm \(\ln x\) enthält.
\end{definition}

\(Y_\nu(x)\) ist singulär bei \(x = 0\); in physikalischen Problemen, die den Ursprung
einschließen, behält man nur \(J_\nu\).

\subsection{Modifizierte Bessel-Funktionen}

Die Substitution \(x \mapsto ix\) in~\eqref{eq:bessel_de} ergibt die modifizierte
Bessel-Gleichung \(x^2 y'' + xy' - (x^2+\nu^2)y = 0\) mit Lösungen
\[
  I_\nu(x) = i^{-\nu}J_\nu(ix), \qquad
  K_\nu(x) = \frac{\pi}{2}\,\frac{I_{-\nu}(x)-I_\nu(x)}{\sin(\nu\pi)}.
\]
\(K_\nu(x)\) klingt exponentiell ab: \(K_\nu(x) \sim \sqrt{\pi/(2x)}\,e^{-x}\) für
\(x \to +\infty\).

\subsection{Rekurrenzrelationen}

\begin{theorem}[Drei-Term-Rekurrenz]\label{thm:bessel_recur_de}
Für \(x \neq 0\) gelten:
\begin{align}
  J_{\nu-1}(x) + J_{\nu+1}(x) &= \frac{2\nu}{x}\,J_\nu(x),\label{eq:bessel_rec1_de}\\
  J_{\nu-1}(x) - J_{\nu+1}(x) &= 2J_\nu'(x).\label{eq:bessel_rec2_de}
\end{align}
\end{theorem}

\subsection{Wronski-Determinante}

\begin{theorem}[Wronskian von \(J_\nu\) und \(Y_\nu\)]
\[
  W\bigl[J_\nu(x), Y_\nu(x)\bigr]
  = J_\nu(x)Y_\nu'(x) - J_\nu'(x)Y_\nu(x)
  = \frac{2}{\pi x}.
\]
\end{theorem}

Die nicht-verschwindende Wronski-Determinante bestätigt, dass \(J_\nu\) und \(Y_\nu\)
ein Fundamentalsystem für \(x > 0\) bilden.

\subsection{Nullstellen und asymptotisches Verhalten}

Alle Nullstellen von \(J_\nu\) für \(\nu > -1\) sind reell und einfach.
Asymptotisch gilt für \(x \to +\infty\):
\[
  J_\nu(x) \sim \sqrt{\frac{2}{\pi x}}\cos\!\left(x - \frac{\nu\pi}{2} - \frac{\pi}{4}\right).
\]

%% ============================================================
\section{Hypergeometrische Funktionen}
\label{sec:hypergeometric}
%% ============================================================

\subsection{Die Gaußsche hypergeometrische Reihe}

\begin{definition}[Pochhammer-Symbol (aufsteigende Fakultät)]
\[
  (a)_0 = 1, \qquad (a)_n = a(a+1)\cdots(a+n-1) = \frac{\Gamma(a+n)}{\Gamma(a)}.
\]
\end{definition}

\begin{definition}[Gaußsche hypergeometrische Funktion]\label{def:2F1_de}
Für \(c \notin \{0,-1,-2,\ldots\}\) und \(|z| < 1\) sei
\[
  {}_2F_1(a,b;c;z)
  = \sum_{n=0}^\infty \frac{(a)_n(b)_n}{(c)_n\,n!}\,z^n.
\]
\end{definition}

Die Reihe konvergiert absolut für \(|z| < 1\) und kann analytisch auf
\(\mathbb{C}\setminus[1,\infty)\) fortgesetzt werden.

\subsection{Gaußsche hypergeometrische Differentialgleichung}

\begin{theorem}
\({}_2F_1(a,b;c;z)\) erfüllt die Gleichung
\[
  z(1-z)w'' + \bigl[c-(a+b+1)z\bigr]w' - ab\,w = 0.
\]
Diese Gleichung besitzt reguläre singuläre Punkte bei \(z = 0, 1, \infty\).
\end{theorem}

\subsection{Transformationsformeln}

\begin{theorem}[Pfaffsche Transformation]
\[
  {}_2F_1(a,b;c;z) = (1-z)^{-a}\,{}_2F_1\!\left(a,c-b;c;\frac{z}{z-1}\right).
\]
\end{theorem}

\begin{theorem}[Eulersche Transformation]
\[
  {}_2F_1(a,b;c;z) = (1-z)^{c-a-b}\,{}_2F_1(c-a,c-b;c;z).
\]
\end{theorem}

\begin{theorem}[Kummers 24 Lösungen]
Kummer fand 24 Lösungen der hypergeometrischen Gleichung, die durch Anwendung der sechs
Symmetrien der Dreiergruppe \(\{0,1,\infty\}\) und der zwei lokalen Lösungen an jedem
Punkt entstehen.
\end{theorem}

\subsection{Wichtige Spezialfälle}

\begin{proposition}[Sonderfälle]
\begin{align*}
  {}_2F_1(1,1;1;z) &= \frac{1}{1-z} \quad (|z|<1),\\
  {}_2F_1\!\left(\tfrac{1}{2},\tfrac{1}{2};1;k^2\right) &= \frac{2}{\pi}\,K(k),\\
  P_n(x) &= {}_2F_1\!\left(-n,n+1;1;\tfrac{1-x}{2}\right),\\
  \ln(1+z) &= z\,{}_2F_1(1,1;2;-z).
\end{align*}
\end{proposition}

\subsection{Konfluente hypergeometrische (Kummer-) Funktionen}

Für \(b \to \infty\) mit \(z \mapsto z/b\) in Definition~\ref{def:2F1_de} entsteht die
konfluente hypergeometrische Funktion (Kummer-Funktion):
\[
  M(a;c;z) = {}_1F_1(a;c;z)
  = \sum_{n=0}^\infty \frac{(a)_n}{(c)_n\,n!}\,z^n,
\]
welche die Kummer-Gleichung \(z w'' + (c-z)w' - aw = 0\) erfüllt.
Bemerkenswerterweise gilt \(M(a;a;z) = e^z\) und
\(\operatorname{erf}(z) = (2z/\sqrt{\pi})\,M(1/2;3/2;-z^2)\).

%% ============================================================
\section{Legendre-Polynome}
\label{sec:legendre}
%% ============================================================

\subsection{Definition via Rodrigues-Formel}

\begin{theorem}[Rodrigues-Formel]\label{thm:rodrigues_de}
\[
  P_n(x) = \frac{1}{2^n n!}\,\frac{d^n}{dx^n}(x^2-1)^n.
\]
\end{theorem}

Die ersten Polynome lauten:
\[
  P_0(x)=1,\quad P_1(x)=x,\quad P_2(x)=\tfrac{1}{2}(3x^2-1),\quad
  P_3(x)=\tfrac{1}{2}(5x^3-3x).
\]

\subsection{Orthogonalität}

\begin{theorem}[Orthogonalität der Legendre-Polynome]\label{thm:legendre_ortho_de}
\[
  \int_{-1}^1 P_n(x)\,P_m(x)\,dx = \frac{2}{2n+1}\,\delta_{nm}.
\]
\end{theorem}

\begin{proof}
Beide Polynome \(P_n\) und \(P_m\) genügen der Legendre-Gleichung.
Man multipliziert die Gleichung für \(P_n\) mit \(P_m\), subtrahiert das analoge
Produkt für \(P_m\) mal \(P_n\) und integriert über \([-1,1]\).
Die Randterme verschwinden, da \((1-x^2)\) bei \(\pm 1\) null ist, was
\([n(n+1)-m(m+1)]\int_{-1}^1 P_n P_m\,dx = 0\) ergibt.
Für \(n \neq m\) folgt daher, dass das Integral null ist.
Die Normierung \(2/(2n+1)\) ergibt sich aus \(\|P_n\|^2\) via Rodrigues-Formel
und \(n\)-facher partieller Integration.
\end{proof}

\subsection{Erzeugende Funktion}

\begin{theorem}[Erzeugende Funktion]
Für \(|t| < 1\) und \(x \in [-1,1]\) gilt
\[
  \frac{1}{\sqrt{1-2xt+t^2}} = \sum_{n=0}^\infty P_n(x)\,t^n.
\]
\end{theorem}

\subsection{Zugeordnete Legendre-Funktionen und Kugelflächenfunktionen}

\begin{definition}[Zugeordnete Legendre-Funktion]
Für \(0 \leq m \leq n\) sei
\[
  P_n^m(x) = (-1)^m(1-x^2)^{m/2}\frac{d^m}{dx^m}P_n(x).
\]
\end{definition}

\begin{definition}[Kugelflächenfunktionen]
\[
  Y_n^m(\theta,\varphi)
  = \sqrt{\frac{(2n+1)}{4\pi}\frac{(n-|m|)!}{(n+|m|)!}}\,
  P_n^{|m|}(\cos\theta)\,e^{im\varphi},
\]
für \(n \geq 0\), \(-n \leq m \leq n\).  Die Kugelflächenfunktionen bilden eine
vollständige Orthonormalbasis von \(L^2(S^2)\).
\end{definition}

%% ============================================================
\section{Hermite- und Laguerre-Polynome}
\label{sec:hermite_laguerre}
%% ============================================================

\subsection{Hermite-Polynome}

\begin{definition}[Hermite-Polynome, Physiker-Konvention]
\[
  H_n(x) = (-1)^n e^{x^2}\frac{d^n}{dx^n}e^{-x^2}.
\]
\end{definition}

Die ersten Polynome sind \(H_0 = 1\), \(H_1 = 2x\), \(H_2 = 4x^2-2\), \(H_3 = 8x^3-12x\).

\begin{theorem}[Orthogonalität]
\[
  \int_{-\infty}^\infty H_n(x)H_m(x)e^{-x^2}\,dx = 2^n\,n!\,\sqrt{\pi}\,\delta_{nm}.
\]
\end{theorem}

\begin{theorem}[Rekurrenzrelation]
\[
  H_{n+1}(x) = 2x\,H_n(x) - 2n\,H_{n-1}(x).
\]
\end{theorem}

\subsection{Quantenmechanischer harmonischer Oszillator}

Die stationäre Schrödinger-Gleichung für den harmonischen Oszillator
\[
  -\frac{\hbar^2}{2m}\psi'' + \frac{1}{2}m\omega^2 x^2\psi = E\psi
\]
vereinfacht sich mit \(\xi = \sqrt{m\omega/\hbar}\,x\) zu
\[
  \psi'' - (\xi^2 - \lambda)\psi = 0, \qquad \lambda = 2E/(\hbar\omega).
\]
Quadratintegrierbare Lösungen existieren nur für \(\lambda = 2n+1\), \(n \in \mathbb{N}_0\),
was quantisierte Energien \(E_n = \hbar\omega(n+\tfrac{1}{2})\) und Wellenfunktionen
\[
  \psi_n(x) = \left(\frac{m\omega}{\pi\hbar}\right)^{1/4}
  \frac{1}{\sqrt{2^n n!}}\,H_n\!\left(\sqrt{\frac{m\omega}{\hbar}}\,x\right)
  e^{-m\omega x^2/(2\hbar)}
\]
ergibt.

\subsection{Laguerre-Polynome}

\begin{definition}[Laguerre-Polynome]
\[
  L_n(x) = \frac{e^x}{n!}\frac{d^n}{dx^n}(x^n e^{-x}) = \sum_{k=0}^n \binom{n}{k}
  \frac{(-x)^k}{k!}.
\]
\end{definition}

\begin{theorem}[Orthogonalität]
\[
  \int_0^\infty L_n(x)L_m(x)e^{-x}\,dx = \delta_{nm}.
\]
\end{theorem}

\begin{theorem}[Rekurrenzrelation]
\[
  (n+1)L_{n+1}(x) = (2n+1-x)L_n(x) - n\,L_{n-1}(x).
\]
\end{theorem}

\subsection{Wasserstoffatom}

Die radiale Schrödinger-Gleichung für das wasserstoffähnliche Atom führt nach
Variablentrennung auf die zugeordnete Laguerre-Gleichung.
Die Radialwellenfunktionen lauten
\[
  R_{nl}(r) = -\sqrt{\left(\frac{2}{na_0}\right)^3\frac{(n-l-1)!}{2n[(n+l)!]^3}}\,
  e^{-r/(na_0)}\!\left(\frac{2r}{na_0}\right)^l L_{n-l-1}^{2l+1}\!\!\left(\frac{2r}{na_0}\right),
\]
wobei \(a_0\) der Bohrsche Radius und \(L_n^\alpha\) die verallgemeinerten Laguerre-Polynome
\(L_n^\alpha(x) = e^x x^{-\alpha}(d^n/dx^n)(e^{-x}x^{n+\alpha})/n!\) sind.

%% ============================================================
\section{Elliptische Funktionen}
\label{sec:elliptic}
%% ============================================================

\subsection{Elliptische Integrale}

\begin{definition}[Vollständige elliptische Integrale]
Für den Modulus \(0 \leq k < 1\) sei
\begin{align*}
  K(k) &= \int_0^{\pi/2}\frac{d\theta}{\sqrt{1-k^2\sin^2\theta}}, \\
  E(k) &= \int_0^{\pi/2}\sqrt{1-k^2\sin^2\theta}\,d\theta.
\end{align*}
\end{definition}

\begin{theorem}[Legendre-Relation]
Sei \(k' = \sqrt{1-k^2}\) der komplementäre Modulus.  Dann gilt
\[
  K(k)\,E(k') + K(k')\,E(k) - K(k)\,K(k') = \frac{\pi}{2}.
\]
\end{theorem}

\subsection{Weierstraßsche elliptische Funktion}

\begin{definition}[Weierstraßsche \(\wp\)-Funktion]
Sei \(\Lambda = \omega_1\mathbb{Z} + \omega_2\mathbb{Z}\) ein Gitter in \(\mathbb{C}\).
\[
  \wp(z; \Lambda) = \frac{1}{z^2}
  + \sum_{\omega \in \Lambda\setminus\{0\}} \left(\frac{1}{(z-\omega)^2} - \frac{1}{\omega^2}\right).
\]
\end{definition}

\begin{theorem}[Differentialgleichung für \(\wp\)]
\[
  (\wp')^2 = 4\wp^3 - g_2\,\wp - g_3,
\]
wobei die Eisenstein-Reihen \(g_2 = 60 G_4(\Lambda)\) und \(g_3 = 140 G_6(\Lambda)\).
Die Diskriminante \(\Delta = g_2^3 - 27g_3^2 \neq 0\) stellt sicher, dass das kubische
Polynom rechts voneinander verschiedene Wurzeln hat.
\end{theorem}

\begin{remark}
Die Gleichung \(y^2 = 4x^3 - g_2 x - g_3\) ist die Weierstraß-Normalform einer
elliptischen Kurve.  Jede elliptische Kurve über \(\mathbb{C}\) ist isomorph zu
\(\mathbb{C}/\Lambda\) für ein geeignetes Gitter \(\Lambda\), und die Parametrisierung
\(z \mapsto (\wp(z), \wp'(z))\) realisiert diesen Isomorphismus explizit.
\end{remark}

\subsection{Jacobische elliptische Funktionen}

Die Jacobischen Funktionen \(\operatorname{sn}(u,k)\), \(\operatorname{cn}(u,k)\),
\(\operatorname{dn}(u,k)\) sind als Umkehrungen des elliptischen Integrals erster Art
definiert.  Sie erfüllen
\begin{align*}
  \operatorname{sn}^2 + \operatorname{cn}^2 &= 1, \\
  k^2\operatorname{sn}^2 + \operatorname{dn}^2 &= 1, \\
  \frac{d}{du}\operatorname{sn}(u,k) &= \operatorname{cn}(u,k)\,\operatorname{dn}(u,k).
\end{align*}
Die Perioden sind \(4K(k)\) für \(\operatorname{sn}\) und \(\operatorname{cn}\) sowie
\(2K(k)\) für \(\operatorname{dn}\).

%% ============================================================
\section{Riemannsche Zetafunktion und Dirichlet-\texorpdfstring{$L$}{L}-Funktionen}
\label{sec:zeta}
%% ============================================================

\subsection{Die Riemannsche Zetafunktion}

\begin{definition}[Riemannsche Zetafunktion]
Für \(\Re(s) > 1\) sei
\[
  \zeta(s) = \sum_{n=1}^\infty \frac{1}{n^s} = \prod_p \frac{1}{1-p^{-s}},
\]
das Produkt über alle rationalen Primzahlen \(p\).
\end{definition}

\subsection{Analytische Fortsetzung und Funktionalgleichung}

\begin{theorem}[Funktionalgleichung von \(\zeta\)]\label{thm:zeta_functional_de}
Die Funktion \(\zeta(s)\) setzt sich zu einer meromorphen Funktion auf ganz \(\mathbb{C}\)
fort mit einem einfachen Pol bei \(s = 1\) mit Residuum \(1\).
Die vervollständigte Zetafunktion
\[
  \xi(s) = \tfrac{1}{2}s(s-1)\pi^{-s/2}\Gamma\!\left(\tfrac{s}{2}\right)\zeta(s)
\]
erfüllt die Funktionalgleichung
\[
  \xi(s) = \xi(1-s).
\]
Äquivalent dazu ist
\[
  \zeta(s) = 2^s\pi^{s-1}\sin\!\left(\frac{\pi s}{2}\right)\Gamma(1-s)\,\zeta(1-s).
\]
\end{theorem}

\begin{proof}[Beweisidee]
Ausgangspunkt ist die Mellin-Transformation \(\pi^{-s/2}\Gamma(s/2)\zeta(s)
= \int_0^\infty x^{s/2-1}\theta(x)\,dx\) mit
\(\theta(x) = \sum_{n=1}^\infty e^{-\pi n^2 x}\).
Die modulare Transformationsformel
\(\theta(1/x) = x^{1/2}\theta(x) + (x^{1/2}-1)/2\)
(Konsequenz der Jacobischen Theta-Identität) liefert nach Aufteilen des Integrals
bei \(x = 1\) die Symmetrie \(s \leftrightarrow 1-s\).
\end{proof}

\subsection{Spezielle Werte}

\begin{theorem}[Eulersche Formel für gerade Argumente]
Für \(n \in \mathbb{N}\) gilt
\[
  \zeta(2n) = \frac{(-1)^{n+1}(2\pi)^{2n}B_{2n}}{2\,(2n)!},
\]
wobei \(B_{2n}\) die Bernoulli-Zahlen sind.  Insbesondere:
\(\zeta(2) = \pi^2/6\), \(\zeta(4) = \pi^4/90\), \(\zeta(6) = \pi^6/945\).
\end{theorem}

\subsection{Dirichlet-\(L\)-Funktionen}

\begin{definition}[Dirichlet-Charakter und \(L\)-Funktion]
Ein Dirichlet-Charakter modulo \(q\) ist eine vollständig multiplikative Funktion
\(\chi: \mathbb{Z} \to \mathbb{C}\) mit \(\chi(n+q) = \chi(n)\) und \(\chi(n) = 0\)
falls \(\gcd(n,q) > 1\).  Die zugehörige \(L\)-Funktion ist
\[
  L(s,\chi) = \sum_{n=1}^\infty \frac{\chi(n)}{n^s}
  = \prod_p \frac{1}{1-\chi(p)p^{-s}}, \qquad \Re(s) > 1.
\]
\end{definition}

\begin{theorem}[Nicht-Verschwinden bei \(s=1\)]
Für jeden nicht-Hauptcharakter \(\chi\) gilt \(L(1,\chi) \neq 0\).
Dies ist die zentrale analytische Zutat für Dirichlets Satz über Primzahlen in
arithmetischen Progressionen.
\end{theorem}

\begin{theorem}[Funktionalgleichung für \(L(s,\chi)\)]
Sei \(\chi\) ein primitiver Charakter modulo \(q\) mit Wurzelzahl \(\varepsilon(\chi)\).
Definiere
\[
  \Lambda(s,\chi) = \left(\frac{q}{\pi}\right)^{(s+a)/2}\!\Gamma\!\left(\frac{s+a}{2}\right)L(s,\chi),
\]
wobei \(a = 0\) falls \(\chi(-1) = 1\) und \(a = 1\) falls \(\chi(-1) = -1\).  Dann gilt
\[
  \Lambda(s,\chi) = \frac{\tau(\chi)}{i^a\sqrt{q}}\,\Lambda(1-s,\bar\chi),
\]
wobei \(\tau(\chi) = \sum_{n=1}^q \chi(n)e^{2\pi in/q}\) die Gaußsche Summe ist.
\end{theorem}

%% ============================================================
\section{Anwendungen}
\label{sec:anwendungen}
%% ============================================================

\subsection{Quantenmechanik}

Spezielle Funktionen durchdringen die gesamte Quantenmechanik.
Die zeitunabhängige Schrödinger-Gleichung in kugelsymmetrischen Potentialen
trennt sich in einen Radial- und einen Winkelanteil auf:
\begin{itemize}
  \item \textbf{Winkelanteil}: Kugelflächenfunktionen \(Y_n^m(\theta,\varphi)\)
        (Legendre-Polynome + Exponentialfunktionen).
  \item \textbf{Harmonischer Oszillator}: Hermite-Polynome \(H_n\).
  \item \textbf{Wasserstoffatom}: zugeordnete Laguerre-Polynome \(L_{n-l-1}^{2l+1}\).
  \item \textbf{Lineares Potential / WKB}: Airy-Funktionen \(\operatorname{Ai}(x)\),
        \(\operatorname{Bi}(x)\).
  \item \textbf{Streuung}: Bessel- und sphärische Bessel-Funktionen \(j_l(kr)\).
\end{itemize}

\subsection{Klassische Elektrodynamik}

Die Laplace-Gleichung \(\nabla^2\Phi = 0\) in Zylinderkoordinaten besitzt Lösungen in
Form von Bessel-Funktionen \(J_n(kr)\) und \(Y_n(kr)\).  In Kugelkoordinaten treten
Kugelflächenfunktionen auf.

\subsection{Zahlentheorie und Primzahlverteilung}

Der Primzahlsatz besagt \(\pi(x) \sim x/\ln x\) für \(x \to \infty\).
Sein Beweis hängt am Nicht-Verschwinden von \(\zeta(s)\) auf der Geraden \(\Re(s) = 1\)
und an der expliziten Formel
\[
  \psi(x) = x - \sum_\rho \frac{x^\rho}{\rho} - \ln(2\pi) - \tfrac{1}{2}\ln(1-x^{-2}),
\]
wobei die Summe über die nichttrivialen Nullstellen \(\rho\) von \(\zeta\) läuft.

\begin{conjecture}[Riemannsche Hypothese]
Alle nichttrivialen Nullstellen von \(\zeta(s)\) erfüllen \(\Re(\rho) = 1/2\).
Diese Vermutung ist eines der sieben Clay-Millennium-Probleme und bis heute unbewiesen.
\end{conjecture}

\subsection{Elliptische Funktionen und Zahlentheorie}

Die \(j\)-Invariante
\[
  j(\tau) = 1728\,\frac{g_2^3}{\Delta} = q^{-1} + 744 + 196884\,q + \cdots,
  \quad q = e^{2\pi i\tau},
\]
ist der Hauptmodul für die Modulgruppe \(\mathrm{SL}(2,\mathbb{Z})\) und spielt eine
zentrale Rolle in der Theorie der komplexen Multiplikation und im Moonshine-Phänomen.

%% ============================================================
\section{Literatur}
%% ============================================================

\begin{thebibliography}{99}

\bibitem{AS}
M.~Abramowitz und I.~A.~Stegun,
\textit{Handbook of Mathematical Functions with Formulas, Graphs, and Mathematical Tables},
Dover Publications, New York, 1972.

\bibitem{NIST}
F.~W.~J.~Olver, D.~W.~Lozier, R.~F.~Boisvert und C.~W.~Clark (Hrsg.),
\textit{NIST Digital Library of Mathematical Functions},
Release 1.2.0, \url{https://dlmf.nist.gov/}, 2024.

\bibitem{WhittakerWatson}
E.~T.~Whittaker und G.~N.~Watson,
\textit{A Course of Modern Analysis}, 5.~Aufl.,
Cambridge University Press, 2021 (Nachdruck der Aufl.\ von 1927).

\bibitem{Lebedev}
N.~N.~Lebedev,
\textit{Spezielle Funktionen und ihre Anwendungen} (russ.\ Orig.\ 1963),
deutsche Ausgabe: Bibliographisches Institut, 1973.

\bibitem{Stein}
E.~M.~Stein und R.~Shakarchi,
\textit{Complex Analysis},
Princeton University Press, 2003.

\bibitem{GradshteynRyzhik}
I.~S.~Gradshteyn und I.~M.~Ryzhik,
\textit{Table of Integrals, Series, and Products}, 8.~Aufl.,
Academic Press, 2015.

\bibitem{Andrews}
G.~E.~Andrews, R.~Askey und R.~Roy,
\textit{Special Functions},
Cambridge University Press, 2000.

\bibitem{Watson}
G.~N.~Watson,
\textit{A Treatise on the Theory of Bessel Functions}, 2.~Aufl.,
Cambridge University Press, 1944.

\bibitem{Silverman}
J.~H.~Silverman,
\textit{The Arithmetic of Elliptic Curves}, 2.~Aufl.,
Springer, 2009.

\end{thebibliography}

\end{document}
